\section{Introduction}\label{sec:intro}

\subsection{Background}

Line and pseudoline arrangements with many triangular faces have been studied by several authors. In the affine setting this is the Kobon triangle problem~\cite{fujimura78,a006066}, which asks for the maximum number of triangular faces determined by $n$ lines in the plane. A brief survey of the literature and the basic notions is given in our previous paper~\cite{pu26b}, which deals with enumeration and classification algorithms for triangle-maximal affine arrangements of an odd number of pseudolines. There we also proved that the enumeration algorithms miss no arrangements, obtained the exhaustive enumeration for $n=27$ pseudolines, and found examples up to $n=93$.

The even case turned out to be harder: its arrangements have a more complicated structure, and, unlike in the odd case, no known upper bound for arbitrary arrangements matches the known examples for any $n\ge6$. Consider the bounds first. Roudneff~\cite{roudneff96} proved that the number $p_3$ of triangles in a projective arrangement $\mathcal{A}$ of $N$ pseudolines satisfies
\begin{equation}\label{eq:p3-bound}
p_3(\mathcal{A})\le\left\lfloor\frac{N(N-1)}{3}\right\rfloor,\qquad N\ge 9.
\end{equation}

Felsner and Kriegel claim that the proof can be adapted to the affine case~\cite[Theorem~2(3)]{fk99}. This gives a bound on the number $a_3$ of triangles in an affine arrangement $\mathcal{A}$ of $n$ pseudolines:
\begin{equation}\label{eq:a3-bound}
a_3(\mathcal{A})\le\left\lfloor\frac{n(n-2)}{3}\right\rfloor.
\end{equation}

Both bounds are stated for arbitrary pseudoline arrangements. For \emph{simple} arrangements, which have no point where three or more pseudolines cross (\emph{multiple points}), \eqref{eq:a3-bound} follows from an elementary segment count~\cite[Subsection~2.3]{pu26b}. The other $n-1$ pseudolines cut each pseudoline into $n-2$ bounded segments, and each segment is a side of at most one triangle. Let $u(\mathcal A)$ be the number of \emph{unused segments}, which are sides of no triangle. Counting all segments gives
\begin{equation}\label{eq:simple-segm}
3a_3(\mathcal{A})+u(\mathcal{A})=n(n-2),
\end{equation}
which implies \eqref{eq:a3-bound}.

For odd $n$, some simple arrangements have no unused segments and attain \eqref{eq:a3-bound} without rounding. Such arrangements are called \emph{perfect}.
By divisibility, they can exist only for $n\equiv 3,5\pmod 6$.
Their triangles alternate along every pseudoline, and the alternation is consistent across the whole arrangement: in a checkerboard coloring, all triangles are black and all other faces are white~\cite[Lemma~2.2]{pu26b}.

In a simple arrangement with even $n$, the alternation cannot propagate: for every pseudoline $\ell$, some other pseudoline has an unused segment with an endpoint on $\ell$~\cite[Proposition~2.3]{blanc11}. Each such segment serves at most two pseudolines, so $u\ge n/2$. This gives Blanc's bound~\cite[Corollary~2.4]{blanc11}:
\begin{equation}\label{eq:blanc-even}
a_{3}(\mathcal{A}_\text{simple})\le\left\lfloor\dfrac{n\left(n - \frac52\right)}{3}\right\rfloor=\begin{cases}
\dfrac{n\left(n - \frac52\right)}{3}, & n\equiv0,4\pmod 6 \\
\dfrac{n\left(n - \frac52\right) - 2}{3}. &  n\equiv2\pmod 6
\end{cases}
\end{equation}
Compared with \eqref{eq:a3-bound}, the bound \eqref{eq:blanc-even} loses $n/6$ triangles, up to rounding, to the $n/2$ unavoidable unused segments.

In the even case, however, one can no longer restrict attention to simple arrangements.
Two examples were already known in which multiple points yield an extra triangle.
The arrangement of $n=8$ lines in \cref{fig:seed} has 15 triangles, while \eqref{eq:blanc-even} allows only 14.
For $n=12$, an arrangement found by Kabanovich~\cite{khodulev99} attains \eqref{eq:blanc-even} with 38 triangles, although no simple arrangement has 38 triangles (see \cref{tab:triple-points}).
We found a third example, $n=14$ with 54 triangles, and it was the starting point of this work. Maiorana~\cite{maiorana26} found it independently while this paper was in preparation.
Still, none of these examples attains the general bound \eqref{eq:a3-bound}, and the literature gives no intermediate bound for non-simple arrangements.

Thus in the even case triangles cannot be packed with the unbroken alternation of perfect arrangements: parity or multiple points break it.
The breaks are reflected in the gap between \eqref{eq:a3-bound} and the known examples, which grows with $n$.
They also affect the computer search: in the perfect case a branch can be discarded at the first break of the alternation, whereas in the even case the search has to work with a budget on the number of triangles (cf.~\cite[Section~3]{pu26b}).

A further difficulty of the even case is that the known infinite series of even line arrangements are obtained from odd ones~\cite{for,bbl} by adding a line far from all intersection points~\cite{bbl,blanc11,savchuk25}, and they do not exceed \eqref{eq:blanc-even}. No construction was previously known that takes an even arrangement as input. The known odd series are $2\cdot2^t+1$, $6\cdot2^t+1$, $14\cdot2^t+1$~\cite{bbl} and $18\cdot2^t+1$, which we found earlier~\cite{pu26}. Iterative doubling~\cite{bbl}, on which these series are based, takes a suitable arrangement of an odd number $n$ of lines and produces an arrangement of $2n-1$ lines.

\subsection{Overview of results}

This paper is concerned with maximizing the number $a_{3}(\mathcal{A})$ of triangular faces in the even case, over simple and non-simple affine arrangements of pseudolines and of straight lines.

We adapt iterative doubling to an even number of lines and apply it to an arrangement of $n=8$ lines. The adaptation temporarily removes one line, applies doubling as many times as needed, and puts the line back. This procedure gives an infinite series of arrangements $\mathcal{C}_t$ of $n=6\cdot2^t+2$ lines with one triangle more than \eqref{eq:blanc-even} allows:
\[
a_3(\mathcal{C}_t)=\frac{n\left(n-\frac52\right)+1}{3}
=\frac{n\left(n-\frac52\right)-2}{3}+1.
\]
The series contains arrangements of $n=8$, $14$, $26$, $50,\ldots$ lines with $a_3=15$, $54$, $204$, $792,\ldots$ triangles. This series (\cref{thm:series}, proved in \cref{sec:series}) constitutes the main part of our results.

\Cref{sec:enum} describes a computer search for pseudoline arrangements, possibly with multiple points. Its results include exhaustive enumerations for $n=6, 8, 10, 12$ and a proof that no arrangement of $n=14$ pseudolines has $a_3>54$. This verifies that the maximum numbers of triangles for these $n$ are 7, 15, 25, 38, 54. The result extends to straight lines: these maxima are attained by the line arrangements mentioned above and by other known ones. We also checked that no arrangement of $n=11$ pseudolines has $a_3=33$, simple or not.

\Cref{sec:augment} describes augmentation, which produces arrangements of an even number of pseudolines from odd ones.
Augmenting perfect arrangements of $n-1$ pseudolines gives arrangements that attain \eqref{eq:blanc-even} for $n\equiv0,4\pmod 6$.
For $n\equiv2\pmod 6$, augmentation can exceed \eqref{eq:blanc-even} by one, like the series above. We found such augmentations for $n=8,14,20,26,32,38,44$.

We found realizations by straight lines of all arrangements, up to isomorphism, given by the exhaustive enumerations of \cref{sec:enum} for even $n$, except two arrangements of $n=10$ pseudolines. These two are not stretchable (verified by the method of biquadratic final polynomials~\cite{br90}). Line realizations of augmented arrangements (\cref{sec:augment}) of $n=20$ and $n=32$ pseudolines with 117 and 315 triangles were found as well. All these line arrangements are collected in \cref{app:figures}.

We also applied the doubling of \cref{sec:series} to pseudolines. Starting from the augmented arrangements with $n=20$ and $n=32$ (\cref{fig:n20-collapse,fig:n32}), we obtained the members with $t\le3$ of the pseudoline series $18\cdot2^t+2$ and $30\cdot2^t+2$, up to $n=146$ and $n=242$. Each of them exceeds \eqref{eq:blanc-even} by one. The members with $n=38$ and $n=74$, with 450 and 1764 triangles, are also realizable by straight lines (\cref{fig:n38,tab:n74-lines}).

These results lead us to conjecture that every arrangement of an even number $n$ of pseudolines satisfies
\begin{equation}\label{eq:conj-bound}
a_{3}\le\left\lfloor\dfrac{n(n - 2)}{3}\right\rfloor-\left\lfloor\dfrac{n}{6}\right\rfloor=\begin{cases}
\dfrac{n(n - \frac52)}{3}, & n\equiv0,4\pmod 6 \\
\dfrac{n(n - \frac52) + 1}{3}. &  n\equiv2\pmod 6
\end{cases}
\end{equation}
\Cref{sec:discussion} discusses the results and the conjecture.
